\documentclass[12pt,a4paper]{article}

% --- Kodierung und Sprache ---
\usepackage[utf8]{inputenc}
\usepackage[T1]{fontenc}
\usepackage[ngerman]{babel}
\usepackage{lmodern}

% --- Mathematik ---
\usepackage{amsmath,amssymb,amsfonts,amsthm}
\usepackage{mathtools}

% --- Layout ---
\usepackage[a4paper,margin=2.7cm]{geometry}
\usepackage{setspace}
\usepackage{booktabs}
\usepackage{longtable}
\usepackage{array}
\usepackage{hyperref}
\usepackage{xcolor}

% --- Standardmäßig robuste Grafikeinbindung ---
\usepackage{graphicx}
\DeclareGraphicsExtensions{.pdf,.png,.jpg}
\graphicspath{{./}{fig/}{figs/}{images/}{img/}}
\newcommand{\includegraphicssafe}[2][]{%
	\IfFileExists{#2}{\includegraphics[#1]{#2}}{%
		\fbox{\parbox[c][4cm][c]{0.8\linewidth}{\centering Datei #2 nicht gefunden}}%
	}%
}

% --- Theorem-Umgebungen ---
\newtheorem{satz}{Satz}[section]
\newtheorem{lemma}[satz]{Lemma}
\newtheorem{korollar}[satz]{Korollar}
\newtheorem{definition}[satz]{Definition}
\newtheorem{bemerkung}[satz]{Bemerkung}
\newtheorem{beispiel}[satz]{Beispiel}

% --- Makros ---
\newcommand{\N}{\mathbb{N}}
\newcommand{\Z}{\mathbb{Z}}
\newcommand{\EFam}{\mathcal{E}}
\newcommand{\AFam}{\mathcal{A}}
\newcommand{\BFam}{\mathcal{B}}
\newcommand{\CFam}{\mathcal{C}}\newcommand{\CC}{\mathcal{C}}
\newcommand{\Fam}{\operatorname{Fam}}
\newcommand{\sign}{\sigma}
\newcommand{\Mass}{\operatorname{M}}
\newcommand{\QE}{Q_E}
\newcommand{\QA}{Q_A}
\newcommand{\QB}{Q_B}
\newcommand{\QC}{Q_C}
\newcommand{\QK}{Q_K}
\newcommand{\Qtot}{Q_{\mathrm{tot}}}
\newcommand{\rank}{\operatorname{f}}
\newcommand{\core}{\operatorname{K}}
\newcommand{\shell}{\operatorname{S}}
\newcommand{\PS}{\operatorname{PS}}

\title{\textbf{Ein arithmetisches Zahlenperiodensystem\\
		auf Basis der EABC-Familien}}
\author{Thomas Hoffbauer}
\date{\today}

\begin{document}
	\maketitle
	
	\begin{abstract}
		Wir entwickeln ein rein mathematisches Periodensystem der natürlichen Zahlen auf Basis der vier Primzahlfamilien
		\[
		E=\{p\text{ prim}: p\equiv 1\!\!\pmod{12}\},\quad
		A=\{p\text{ prim}: p\equiv 5\!\!\pmod{12}\},
		\]
		\[
		B=\{p\text{ prim}: p\equiv 7\!\!\pmod{12}\},\quad
		C=\{p\text{ prim}: p\equiv 11\!\!\pmod{12}\}.
		\]
		Jede natürliche Zahl $n>1$ wird durch ihre Familienzerlegung, ihren Familienrang, ihre Endsignatur sowie durch arithmetische Ladungszahlen beschrieben. Daraus entsteht eine hierarchische Tafel mit chemieähnlicher Struktur: Spalten entsprechen Endsignaturen, Zeilen dem Familienrang, und zusätzliche Größen wie Masse, Schalenladung und Kernladung liefern eine Feinstruktur. Echte Primzahlvierlinge erscheinen dabei als vollständig besetzte neutrale Objekte in der Klasse $(4,E)$.
	\end{abstract}
	
	\tableofcontents
	\bigskip
	
	\section{Grundidee}
	
	Das klassische Periodensystem der Chemie ordnet Atome nach Ordnungszahl, Schalenstruktur und Valenzeigenschaften. Wir suchen ein arithmetisches Analogon, in dem natürliche Zahlen nach der Struktur ihrer Primfaktorzerlegung klassifiziert werden.
	
	Die zentrale Beobachtung lautet: Für Primzahlen $p>3$ kommen modulo $12$ nur die vier Restklassen
	\[
	1,\;5,\;7,\;11
	\]
	vor. Diese definieren vier Primfamilien, die in einer natürlichen Weise eine algebraische Viererstruktur tragen.
	
	\section{Die vier Primfamilien}
	
	\begin{definition}[EABC-Familien]
		Für Primzahlen $p>3$ definieren wir:
		\[
		E:=\{p\text{ prim}: p\equiv 1\pmod{12}\},
		\qquad
		A:=\{p\text{ prim}: p\equiv 5\pmod{12}\},
		\]
		\[
		B:=\{p\text{ prim}: p\equiv 7\pmod{12}\},
		\qquad
		C:=\{p\text{ prim}: p\equiv 11\pmod{12}\}.
		\]
	\end{definition}
	
	\begin{bemerkung}
		Die kleinen Primzahlen $2$ und $3$ stehen außerhalb dieses Vierfamilien-Schemas und werden im Folgenden als \emph{Sonderprimzahlen} behandelt. Für das Periodensystem konzentrieren wir uns auf Zahlen, deren Primfaktorzerlegung nur Primzahlen $>3$ enthält.
	\end{bemerkung}
	
	\section{Algebraische Endsignatur}
	
	\begin{definition}[Familiengruppe]
		Wir identifizieren die Familienmenge $\{E,A,B,C\}$ mit der Kleinschen Vierergruppe
		\[
		V_4\cong \Z/2\Z\times \Z/2\Z.
		\]
		Dabei gilt
		\[
		A^2=B^2=C^2=E,
		\qquad
		AB=C,\quad AC=B,\quad BC=A,
		\qquad
		ABC=E.
		\]
	\end{definition}
	
	\begin{definition}[Familienzerlegung]
		Sei $n>1$ mit Primfaktorzerlegung
		\[
		n=\prod_{p^\nu\parallel n} p^\nu.
		\]
		Dann definieren wir
		\[
		n_E:=\prod_{\substack{p^\nu\parallel n\\ p\in E}} p^\nu,\quad
		n_A:=\prod_{\substack{p^\nu\parallel n\\ p\in A}} p^\nu,\quad
		n_B:=\prod_{\substack{p^\nu\parallel n\\ p\in B}} p^\nu,\quad
		n_C:=\prod_{\substack{p^\nu\parallel n\\ p\in C}} p^\nu.
		\]
		Es gilt eindeutig
		\[
		n=n_E\,n_A\,n_B\,n_C.
		\]
	\end{definition}
	
	\begin{definition}[Endsignatur]
		Die \emph{Endsignatur} einer Zahl $n>1$ sei
		\[
		\sign(n)\in\{E,A,B,C\},
		\]
		definiert durch die Parität der Exponenten in den Familien $A,B,C$. Quadratische Beiträge sind signaturmäßig neutral.
	\end{definition}
	
	\begin{satz}[Signatur hängt nur vom quadratfreien Kern ab]
		Ist
		\[
		n=s^2r
		\]
		mit quadratfreiem Kern $r$, so gilt
		\[
		\sign(n)=\sign(r).
		\]
	\end{satz}
	
	\begin{proof}
		Quadratische Faktoren tragen nur gerade Exponenten und verschwinden in der Familiengruppe, da jedes nichttriviale Element Ordnung $2$ hat.
	\end{proof}
	
	\section{Masse und Ladungen}
	
	\begin{definition}[Arithmetische Masse]
		Für $n\in\N_{>0}$ definieren wir die \emph{arithmetische Masse} durch
		\[
		\Mass(n):=n.
		\]
	\end{definition}
	
	\begin{definition}[Familienladungen]
		Für $X\in\{E,A,B,C\}$ sei
		\[
		Q_X(n):=\#\{p\in X: p\mid n\}.
		\]
		Insbesondere heißt
		\[
		\QE(n)
		\]
		die \emph{E-Schalenladung}.
	\end{definition}
	
	\begin{definition}[Kernladung und Gesamtladungszahl]
		Wir definieren
		\[
		\QK(n):=\QA(n)+\QB(n)+\QC(n)
		\]
		als \emph{Kernladung} und
		\[
		\Qtot(n):=\QE(n)+\QA(n)+\QB(n)+\QC(n)
		\]
		als \emph{Gesamtladungszahl}.
	\end{definition}
	
	\begin{bemerkung}
		Die Ladungen zählen \emph{verschiedene} Primzahlen, nicht Exponenten. Sie unterscheiden sich also von der Endsignatur $\sign(n)$.
	\end{bemerkung}
	
	\section{Schalen-Kern-Zerlegung}
	
	\begin{definition}[Schalen- und Kernanteil]
		Wir setzen
		\[
		\shell(n):=n_E,
		\qquad
		\core(n):=n_A n_B n_C.
		\]
		Dann gilt
		\[
		n=\shell(n)\,\core(n).
		\]
	\end{definition}
	
	\begin{satz}[Der E-Anteil ist signaturneutral]
		Für jede Zahl $n>1$ gilt
		\[
		\sign(n)=\sign(\core(n)).
		\]
	\end{satz}
	
	\begin{proof}
		Da $\sign(n_E)=E$ gilt, folgt
		\[
		\sign(n)=\sign(n_E)\sign(\core(n))=E\cdot \sign(\core(n))=\sign(\core(n)).
		\]
	\end{proof}
	
	\begin{korollar}[Signaturerhaltender Aufbau]
		Ist $q\in E$ eine Primzahl, so gilt für alle $n>1$
		\[
		\sign(qn)=\sign(n).
		\]
	\end{korollar}
	
	\begin{proof}
		Sofort aus $\sign(q)=E$.
	\end{proof}
	
	\section{Familienrang und hierarchische Klassen}
	
	\begin{definition}[Aktive Familien und Familienrang]
		Die Menge der aktiven Familien von $n>1$ sei
		\[
		\Fam(n):=\{X\in\{E,A,B,C\}: n_X>1\}.
		\]
		Der \emph{Familienrang} sei
		\[
		\rank(n):=|\Fam(n)|.
		\]
	\end{definition}
	
	\begin{definition}[Hierarchische Klasse]
		Für $r\in\{1,2,3,4\}$ und $\tau\in\{E,A,B,C\}$ sei
		\[
		\mathcal C_{r,\tau}:=\{n>1:\rank(n)=r,\ \sign(n)=\tau\}.
		\]
	\end{definition}
	
	\begin{satz}[Grobe Periodentafel]
		Das arithmetische Periodensystem erster Ordnung besteht aus den Klassen
		\[
		\mathcal C_{r,\tau},
		\qquad
		r\in\{1,2,3,4\},\ \tau\in\{E,A,B,C\}.
		\]
		Es ist vollständig besetzbar.
	\end{satz}
	
	\begin{proof}
		Durch geeignete Produkte von Primzahlen aus den vier Familien lassen sich alle Kombinationen konstruieren. Die Signaturseite folgt aus der Kleinschen Vierergruppe.
	\end{proof}
	
	\section{Das arithmetische Periodensystem}
	
	\begin{definition}[Zahlenperiodensystem]
		Das \emph{arithmetische EABC-Periodensystem} einer Zahl $n>1$ sei das Tupel
		\[
		\PS(n):=\bigl(\rank(n),\sign(n),\QE(n),\QK(n),\Mass(n)\bigr).
		\]
	\end{definition}
	
	\subsection*{Interpretation}
	\begin{itemize}
		\item $\rank(n)$ beschreibt die Zahl aktiver Familien.
		\item $\sign(n)$ ist die Endklasse.
		\item $\QE(n)$ misst die äußere Schalenladung.
		\item $\QK(n)$ misst die Kernladung.
		\item $\Mass(n)$ ist die Ordnungsmasse.
	\end{itemize}
	
	\subsection*{Tabellenachsen}
	\begin{itemize}
		\item \textbf{Spalten:} Endsignatur $E,A,B,C$
		\item \textbf{Zeilen:} Familienrang $1,2,3,4$
		\item \textbf{Feinstruktur:} Ladungen $\QE,\QK$ und Masse $M$
	\end{itemize}
	
	\section{Beispiele der ersten Tafelstufen}
	
	\begin{longtable}{>{\raggedright\arraybackslash}p{2.5cm} >{\raggedright\arraybackslash}p{5.1cm} >{\raggedright\arraybackslash}p{4.2cm} >{\raggedright\arraybackslash}p{2.4cm}}
		\toprule
		\textbf{Zahl} & \textbf{Faktorisierung / Familien} & \textbf{Signatur und Ladungen} & \textbf{Klasse}\\
		\midrule
		\endhead
		
		$13$ & $13\in E$ & $\sign=E,\ \QE=1,\ \QK=0$ & $\mathcal C_{1,E}$\\
		$5$ & $5\in A$ & $\sign=A,\ \QE=0,\ \QK=1$ & $\mathcal C_{1,A}$\\
		$7$ & $7\in B$ & $\sign=B,\ \QE=0,\ \QK=1$ & $\mathcal C_{1,B}$\\
		$11$ & $11\in C$ & $\sign=C,\ \QE=0,\ \QK=1$ & $\mathcal C_{1,C}$\\
		$65$ & $5\cdot 13\in A\cdot E$ & $\sign=A,\ \QE=1,\ \QK=1$ & $\mathcal C_{2,A}$\\
		$35$ & $5\cdot 7\in A\cdot B$ & $\sign=C,\ \QE=0,\ \QK=2$ & $\mathcal C_{2,C}$\\
		$385$ & $5\cdot 7\cdot 11\in A\cdot B\cdot C$ & $\sign=E,\ \QE=0,\ \QK=3$ & $\mathcal C_{3,E}$\\
		$5005$ & $13\cdot 5\cdot 7\cdot 11\in E\cdot A\cdot B\cdot C$ & $\sign=E,\ \QE=1,\ \QK=3$ & $\mathcal C_{4,E}$\\
		\bottomrule
	\end{longtable}
	
	\section{Echte Primzahlvierlinge als vollständig besetzte neutrale Objekte}
	
	\begin{definition}[Echter Primzahlvierling]
		Ein echter Primzahlvierling ist ein Tupel
		\[
		Q(p)=(p,p+2,p+6,p+8),
		\]
		wobei alle vier Zahlen prim sind.
	\end{definition}
	
	\begin{satz}
		Jeder echte Primzahlvierling gehört zur Klasse
		\[
		\mathcal C_{4,E}.
		\]
	\end{satz}
	
	\begin{proof}
		Ein echter Primzahlvierling enthält genau eine Primzahl aus jeder Familie in einer der beiden zyklischen Ordnungen
		\[
		(A,B,C,E)\quad\text{oder}\quad(C,E,A,B).
		\]
		Daher ist der Familienrang $4$ und
		\[
		\sign(Q)=E\cdot A\cdot B\cdot C=ABC=E.
		\]
	\end{proof}
	
	\begin{bemerkung}
		In der Chemieanalogie erscheinen echte Primzahlvierlinge damit als \emph{voll besetzte neutrale Vierfamilienobjekte}.
	\end{bemerkung}
	
	\section{Aufbauprinzip und Periodizität}
	
	\begin{satz}[Neutrale Anregung]
		Sei $q\in E$ prim. Dann gilt:
		\[
		\rank(qn)=
		\begin{cases}
			\rank(n)+1, & \text{falls }E\notin \Fam(n),\\
			\rank(n), & \text{falls }E\in \Fam(n),
		\end{cases}
		\]
		und zugleich
		\[
		\sign(qn)=\sign(n).
		\]
	\end{satz}
	
	\begin{proof}
		Der Rang steigt genau dann, wenn die Familie $E$ neu hinzukommt. Die Signatur bleibt erhalten, weil $q$ neutral ist.
	\end{proof}
	
	\begin{korollar}[Schalenaufbau]
		Jede Zahl aus einer Klasse $\mathcal C_{3,\tau}$ ohne $E$-Anteil kann durch Multiplikation mit einer $E$-Primzahl in eine Zahl aus $\mathcal C_{4,\tau}$ überführt werden.
	\end{korollar}
	
	\section{Eine endliche Tafel bis uranähnliche Zahlen}
	
	Wir wählen nun als endlichen Abschlussbereich die Zahlen
	\[
	1<n\le 92.
	\]
	Dies ist rein arithmetisch eine erste schwere Zone mit einem natürlichen Endpunkt bei der Ordnungsmasse $92$, analog zur chemischen Zahl $92$.
	
	\subsection*{Repräsentanten bis 92}
	Im Folgenden geben wir nicht alle Zahlen bis 92 an, sondern charakteristische Repräsentanten:
	
	\begin{longtable}{>{\raggedright\arraybackslash}p{1.5cm} >{\raggedright\arraybackslash}p{3.2cm} >{\raggedright\arraybackslash}p{4.4cm} >{\raggedright\arraybackslash}p{4.7cm}}
		\toprule
		\textbf{$n$} & \textbf{Familienbild} & \textbf{Zustand} & \textbf{Bemerkung}\\
		\midrule
		\endhead
		
		$5$ & $A$ & $(1,A,0,1,5)$ & einfaches A-Objekt\\
		$7$ & $B$ & $(1,B,0,1,7)$ & einfaches B-Objekt\\
		$11$ & $C$ & $(1,C,0,1,11)$ & einfaches C-Objekt\\
		$13$ & $E$ & $(1,E,1,0,13)$ & einfaches E-Schalenobjekt\\
		$35$ & $AB$ & $(2,C,0,2,35)$ & Zweifamilienbindung\\
		$55$ & $AC$ & $(2,B,0,2,55)$ & Zweifamilienbindung\\
		$65$ & $EA$ & $(2,A,1,1,65)$ & Schale + Kern\\
		$77$ & $BC$ & $(2,A,0,2,77)$ & kernlastige Kopplung\\
		$91$ & $EC$ & $(2,C,1,1,91)$ & schweres Zweifamilienobjekt\\
		\bottomrule
	\end{longtable}
	
	Dabei steht der Zustandsvektor hier in Kurzform für
	\[
	(\rank(n),\sign(n),\QE(n),\QK(n),\Mass(n)).
	\]
	
	\subsection*{Uranähnliche Zone}
	Wir nennen Zahlen mit großer Ordnungsmasse und hoher Familienkomplexität \emph{uranähnlich}, wenn sie
	\[
	\rank(n)=4
	\]
	oder eine große Kernladung besitzen und im schweren Bereich der natürlichen Ordnung liegen.
	
	Dies ist ein rein mathematischer Begriff, der die starke Komplexität einer Zahl im Periodensystem ausdrückt.
	
	\section{Hierarchische Gesamtdeutung}
	
	Das arithmetische Periodensystem besitzt die folgende Hierarchie:
	
	\begin{enumerate}
		\item \textbf{Rang 1:} elementare Einfamilienobjekte
		\item \textbf{Rang 2:} binäre Bindungsobjekte
		\item \textbf{Rang 3:} ternäre Kernobjekte
		\item \textbf{Rang 4:} voll besetzte Vierfamilienobjekte
	\end{enumerate}
	
	Jede dieser Stufen zerfällt weiter in die vier Endsignaturspalten
	\[
	E,\ A,\ B,\ C.
	\]
	
	Die Masse ordnet die Objekte in der natürlichen Abzählung, die Schalenladung beschreibt den äußeren $E$-Anteil, und die Kernladung beschreibt die nichtneutralen Familienanteile.
	
	\section{Ausblick}
	
	Die hier entwickelte Struktur kann in mehrere Richtungen weitergeführt werden:
	
	\begin{itemize}
		\item detaillierte vollständige Tafel aller Zahlen bis zu einer vorgegebenen Ordnungsmasse,
		\item Untersuchung der Dichteverteilung der Klassen $\mathcal C_{r,\tau}$,
		\item Einbau von Exponentenstufen als zusätzlicher ``Isotopenstruktur'',
		\item Einordnung echter Primzahlvierlinge als besondere neutrale Spitzenobjekte,
		\item Erweiterung zu einem darstellbaren visuellen Zahlenperiodensystem.
	\end{itemize}
	
	\begin{bemerkung}
		Die Analogie zur Chemie ist rein strukturell gemeint. Das hier vorgestellte System ist kein physikalisches Atommodell, sondern eine arithmetische Klassifikation natürlicher Zahlen mit schalen- und kernähnlicher Hierarchie.
	\end{bemerkung}
	
\end{document}